%%=============================================================================
%% Samenvatting
%%=============================================================================

% TODO: De "abstract" of samenvatting is een kernachtige (~ 1 blz. voor een
% thesis) synthese van het document.
%
% Een goede abstract biedt een kernachtig antwoord op volgende vragen:
%
% 1. Waarover gaat de bachelorproef?
% 2. Waarom heb je er over geschreven?
% 3. Hoe heb je het onderzoek uitgevoerd?
% 4. Wat waren de resultaten? Wat blijkt uit je onderzoek?
% 5. Wat betekenen je resultaten? Wat is de relevantie voor het werkveld?
%
% Daarom bestaat een abstract uit volgende componenten:
%
% - inleiding + kaderen thema
% - probleemstelling
% - (centrale) onderzoeksvraag
% - onderzoeksdoelstelling
% - methodologie
% - resultaten (beperk tot de belangrijkste, relevant voor de onderzoeksvraag)
% - conclusies, aanbevelingen, beperkingen
%
% LET OP! Een samenvatting is GEEN voorwoord!

%%---------- Nederlandse samenvatting -----------------------------------------
%
% TODO: Als je je bachelorproef in het Engels schrijft, moet je eerst een
% Nederlandse samenvatting invoegen. Haal daarvoor onderstaande code uit
% commentaar.
% Wie zijn bachelorproef in het Nederlands schrijft, kan dit negeren, de inhoud
% wordt niet in het document ingevoegd.

\IfLanguageName{english}{%
\selectlanguage{dutch}
\chapter*{Samenvatting}
\lipsum[1-4]
\selectlanguage{english}
}{}

%%---------- Samenvatting -----------------------------------------------------
% De samenvatting in de hoofdtaal van het document

\chapter*{\IfLanguageName{dutch}{Samenvatting}{Abstract}}

Deze bachelorproef is gericht op de onderzoek naar deployment procedures van branches voor Securex mainframes. Momenteel gebruikt Securex voor haar mainframes Subversion als versiebeheersysteem, maar op dit verouderde systeem blijft onderhoud kostelijk, biedt het weinig opties voor uitbreiding en wordt het belast door een groeiend vraag. Hiervoor heeft Securex plannen om hun mainframes te laten overstappen naar IDz om gebruik te maken van Git branches. De focus ligt op het uitwerken van een branchingstrategie die ondersteuning biedt voor ontwikkeling, testing en acceptatie. Hieruit volgt de hoofdonderzoeksvraag: ``hoe kan Git-branching binnen IDz worden ingezet voor gecontroleerde en traceerbare deployments over meerdere omgevingen?''. De methodologie omvat in de eerste fase de stand van zaken en verkenning van de Securex-infrastructuur. In de tweede fase een onderzoek naar de deploymentprocedures rekening houdend met de noden van Securex. Hier wordt ook een testomgeving of simulatie opgesteld en uitgevoerd. Voor de derde fase wordt de aanpak geëvalueerd op de criteria en zal hieruit conclusies en aanbevelingen gedaan worden voor verdere optimalisatie en mogelijke implementatie voor Securex.
